\section{Threats to Validity}
\label{sec:threats}
This review has several threats to validity that should be stated explicitly.

\textbf{Internal validity.} Although the review follows a Kitchenham-guided protocol and uses PRISMA terminology to report the selection flow, study identification, screening, and extraction still involve researcher judgment, especially for borderline papers in which the SLM component or empirical contribution was not equally central. To reduce this risk, borderline screening decisions were discussed through team consensus, and the synthesis remained tied to the predefined research questions and the same extraction matrix throughout.

The review does not apply a formal quality appraisal instrument to individual primary studies, a step recommended by both PRISMA 2020 and the Kitchenham guidelines~\cite{kitchenham2007guidelines}. This decision reflects the heterogeneity of the corpus: studies span diverse domains, evaluation protocols, and contribution types, making a single quality scale difficult to apply uniformly without introducing subjective scoring variance. Instead, quality-relevant information---reported limitations, baseline availability, and metric completeness---was captured in the extraction matrix (Table~\ref{tab:rq-matrix}) and used to modulate synthesis claims. Specifically, claims are explicitly caveated when the supporting study set lacks a reported baseline or when the primary metric is accuracy alone.

\textbf{Construct validity.} Core constructs such as ``SLM'', ``hybrid architecture'', ``multi-agent system'', and ``efficiency'' are not standardized across the literature. A single hybrid label may refer to token-level verification, retrieval-grounded generation, planner-based decomposition, or multimodal fusion, while efficiency may mean latency, FLOPs, API cost, bandwidth reduction, or energy use. This limits direct comparability because studies grouped under the same high-level category may differ substantially in mechanism. To partially limit this, the review separates broad architecture labels from extracted orchestration mechanisms and discusses recurring mechanism families explicitly. Relatedly, because the architecture-scope eligibility criterion required every included study to already employ some routing, retrieval, verification, fallback, or role-specialization mechanism, the corpus is not suited to testing whether a control layer helps relative to no control layer at all; conclusions about control layers in this review concern their calibration and fit to the task, not their presence versus absence.

\textbf{Conclusion validity.} Cross-study inference is limited by the uneven and inconsistent reporting of evaluation outcomes. Accuracy is reported in most papers, but latency, cost, and energy appear only in a minority, often with different units, hardware assumptions, or deployment settings, which prevents formal meta-analysis and limits claims about efficiency gains across the included studies. Publication and reporting bias may further favor positive hybrid results, so the extraction also recorded reported limitations, and the synthesis emphasizes uneven reporting and task-sensitive trade-offs rather than universal gains.

\textbf{External validity.} The study set is based on a Web of Science-centered search and a bounded review window, which improves transparency but may miss relevant studies from other databases, preprint platforms, or fast-moving venues. Although WoS covers many major peer-reviewed journals and conference proceedings, the field evolves quickly and the current corpus is heavily concentrated in 2025 and in a limited set of application domains. Therefore, the conclusions should be generalized cautiously beyond the covered time period, tasks, and deployment settings.
